\chapter{Information Geometry}

\section{Between Isomorphic and Unrelated}

Chapter 23 closed by admitting what its own machinery could not do. Isomorphism in 𝒲 is a binary verdict, related or not, and most pairs of worlds a reader might actually want to compare fall into neither category cleanly. Two fountains built from slightly different Lagrangian parameters, one a touch more damped than the other, are not isomorphic in Chapter 23's exact sense, and yet calling them simply unrelated discards everything genuinely similar about them. This chapter's task is to replace that binary verdict with a number: not whether two worlds are related, but how far apart they are when they are not identical, and how close, when they are almost the same.

\section{Why Comparing Configurations Directly Does Not Work}

The most obvious approach fails quickly enough to be worth ruling out explicitly. Chapter 18 already equipped the field's configuration space, 𝓜, with a metric g, and it might seem that comparing two worlds is simply a matter of measuring the g-distance between their states. This does not work, for a reason that becomes obvious once stated: g measures distance between two configurations of the same domain X, points within a single world's own configuration space. Two different worlds, in general, are not defined over the same domain at all, and there is no canonical way to line up one courtyard's points with another's before any distance between them can even be asked for. Direct geometric comparison presupposes exactly the correspondence this chapter is trying to establish, and cannot be used to establish it.

\section{Entropy as a Distribution, Worlds as Beliefs}

The more productive approach abandons direct comparison of raw configurations and compares something more abstract instead: what each world believes about itself. Recall \(L_t(x)\), the live possibility set introduced in Chapter 8, the space of outcomes a point has not yet been forced to choose among. At any point where entropy remains nonzero, \(L_t(x)\) is naturally read as inducing a distribution over live outcomes, weighted by how strongly each remains admissible given everything the field currently holds. A world, at any given moment, is then not merely a configuration but a belief, in the technical sense this book borrows from the branch of geometry concerned with spaces of probability distributions: the entirety of its entropy structure, across every point still under active uncertainty, describes a point in a much larger space, the space of possible belief states a persistent field could hold. Comparing two worlds becomes comparing two points in this belief space instead of two points in 𝓜 directly, and unlike raw configurations, belief states do not require the two worlds to share a domain at all, only a common vocabulary of what remains uncertain and by how much.

\section{The Information Metric}

This space of belief states carries its own natural notion of distance, arising not from any external choice but from the distributions themselves: two belief states are close when they are difficult to distinguish from a small number of observations, and far apart when a single observation would reliably tell them apart. The precise quantity this book adopts, without deriving it from scratch, is a Fisher-information-style metric, measuring the local sensitivity of a belief state's distribution to small changes in whatever parameters describe it, and its natural companion, a divergence between two belief states measuring how much is lost, in an information-theoretic sense, by approximating one world's entropy structure with the other's. This divergence is generally asymmetric, treating one world as the reference and the other as the approximation, and this asymmetry is not a defect to be corrected; approximating a richly resolved world with a barely-resolved one loses more than the reverse, and the metric is right to say so. A symmetrized version, averaging the divergence taken in both directions, gives the genuine distance this chapter has been building toward: a single number, computable in principle at any point where both worlds carry active uncertainty, measuring how far apart their beliefs about that point actually are.

\section{From Points to Trajectories}

Consistent with the orbit-primacy this book has insisted on since Chapter 20, a distance defined only at a single moment is not yet a distance between worlds in the sense this chapter needs. Two worlds are properly compared not at one instant but across their whole admissible trajectories, and the natural extension accumulates the pointwise informational distance along corresponding stretches of each world's history: integrated, moment by moment, over as much of each trajectory as can be meaningfully aligned. Two fountains whose instantaneous belief states are nearly identical at every corresponding moment of their respective histories are close in this trajectory-level sense; two fountains that happen to coincide at one instant but diverge sharply before and after are not, however similar that single shared instant might have made them look in isolation. This is the same discipline Chapter 22 insisted on for conservation, applied here to comparison: a claim about one instant is not yet a claim about the orbit, and this chapter's distance, properly stated, is a distance between orbits rather than between snapshots.

\section{Near Misses and Family Resemblance}

This finally answers the question Chapter 23 left open. Two fountains related by no isomorphism in 𝒲, one slightly more damped than the other, now receive a real, computable distance rather than a bare declaration of unrelatedness, and a small distance is precisely what this book means by family resemblance: not sameness of type in Chapter 23's exact sense, but proximity of belief across corresponding trajectories, close enough that mistaking one for the other would cost little in the information-theoretic sense this chapter has made precise. Isomorphism and distance are not competing accounts of similarity. They are complementary: isomorphism asks whether two worlds share an exact architecture; distance asks, when they do not, how expensive the difference actually is.

\section{Closing Part IV}

This chapter completes the ascent Part IV has been climbing since Chapter 18. Chapter 18 asked which states are admissible, and gave the constraint manifold C its geometric home. Chapter 19 asked how admissible states evolve, and gave that evolution an energy functional to descend. Chapter 20 asked what allows persistence without stasis, and gave identity its home as a conserved orbit rather than a static point. Chapter 21 asked how local truths become one coherent world, and gave that question a sheaf-theoretic answer, naming hallucination's spatial failure precisely for the first time. Chapter 22 asked what a genuine history is obligated to preserve, and gave conservation its own trajectory-level meaning, naming hallucination's temporal failure to complete the diagnosis Chapter 21 began. Chapter 23 asked when two different worlds share an architecture, and Chapter 24 has asked, for the much more common case where they do not share one exactly, how far apart they actually are.

Together, these seven chapters have supplied, formally, everything Parts II and III used only informally: C, \(\Pi_C\), argmin, \(F_{\mathrm{phys}}\), fixed points, gluing, conservation, and now comparison. Nothing in the mathematics remains owed. What remains is to build it, and Part V, beginning with Chapter 25, turns from what this framework is to how it actually runs, as software, on real machines, for real persistent worlds.
